% !TeX encoding = UTF-8
\documentclass[variant=docente,cover=institutional,mode=screen]{ua-engineering-v6-article}
% !TeX encoding = UTF-8
\UASetUniversity{Universidad Austral}
\UASetFaculty{Facultad de Ingeniería}
\UASetDepartment{Departamento de Computación}
\UASetCampus{Pilar}
\UASetCareer{Ingeniería Informática}
\UASetCommission{Comisión A}
\UASetProfessor{Equipo docente}
\UASetSemester{Segundo cuatrimestre 2026}
\UASetCourse{Sistemas Distribuidos}
\UASetDate{2026}

\UASetClassNumber{09}
\UASetClassTitle{Arquitecturas distribuidas reales}
\title{Clase 09 --- Guía docente: fronteras, eventos, outbox y sagas}
\author{\UAProfessor}
\date{\UADate}
\begin{document}
\maketitle

\section{Objetivo pedagógico profundo}

La clase debe corregir el sesgo “más distribuido = más avanzado”. El estudiante debe aprender a defender una frontera porque compra autonomía, escalado o aislamiento real, y a reconocer la prima distribuida que paga.

\section{Resultados de aprendizaje}

\begin{itemize}
\item distinguir monolito modular de bola de barro;
\item identificar señales reales para extraer un servicio;
\item detectar monolitos distribuidos;
\item modelar bounded contexts y ownership;
\item elegir RPC, evento, comando o consulta;
\item resolver dual write con outbox;
\item diseñar sagas y compensaciones;
\item proponer migración incremental y capacidad operativa.
\end{itemize}

\section{Arquitectura didáctica v9}

\begin{tabular}{p{0.14\textwidth}p{0.20\textwidth}p{0.57\textwidth}}
\toprule
Tiempo & Bloque & Propósito \\
\midrule
0--15 & Apertura & Plataforma de cursos: qué separar y por qué. \\
15--35 & Monolito modular & Baseline y modularidad antes de distribución. \\
35--60 & Microservicios & Prima distribuida, señales y caso Netflix. \\
60--80 & Dominios/datos & Bounded contexts, ownership y Conway. \\
80--90 & Pausa activa & Clasificar interacciones RPC/evento. \\
90--115 & Eventos/outbox & Comandos, contratos y dual write. \\
115--145 & Sagas & Compensaciones, coreografía y orquestación. \\
145--155 & Pausa & Recuperación breve. \\
155--170 & Migración/operación & Strangler y capacidades previas. \\
170--178 & Taller & Diseñar frontera de la plataforma de cursos. \\
178--180 & Exit ticket & Prima, outbox, compensación. \\
\bottomrule
\end{tabular}

\section{Preparación previa}

\begin{itemize}
\item Llevar un diagrama de monolito modular y el mismo sistema como microservicios.
\item Preparar la cadena de disponibilidad 0.999 elevada a cuatro dependencias.
\item Tener un ejemplo de outbox y una saga con pasos críticos/secundarios.
\item Revisar casos oficiales de Netflix y patrones AWS.
\item Identificar una decisión real del Trabajo Final de cada grupo para discutir.
\end{itemize}

\section{Conducción paso a paso}

\subsection{Apertura}

Pregunte “¿qué separarían primero?” antes de presentar teoría. Registre razones en columnas: dominio, carga, equipo, falla, tecnología. Luego clasifique: solo las primeras cuatro pueden justificar arquitectura; la tecnología por sí sola no.

\subsection{Monolito modular}

Use un baseline fuerte. Si se presenta como opción inferior, los estudiantes no aprenden a medir la prima. Muestre reglas de modularidad y cómo pruebas de arquitectura pueden limitar dependencias.

\subsection{Microservicios}

No enumere ventajas sin costos. Para cada ventaja, pida “¿qué nuevo problema aparece?”. Con Netflix, enfatice escala y ownership, no imitación.

\subsection{Bounded contexts}

Use una palabra ambigua del dominio. Haga que grupos definan modelos distintos y luego contrato de traducción. Esto vuelve concreta la frontera semántica.

\subsection{RPC/eventos}

Entregue cuatro interacciones. Los estudiantes deben elegir y decir qué pasa si el receptor tarda o no existe. Introduzca comando/evento/consulta para mejorar precisión.

\subsection{Dual write y sagas}

Deje que intenten cambiar el orden DB/evento. Cuando vean que no resuelve, introduzca outbox. Para saga, distinga pasos cuyo fallo debe compensar negocio de pasos secundarios que solo se reintentan.

\section{Casos reales y narrativa}

\subsection{Monolith First}

Narrativa: cuando el dominio es incierto, una frontera de proceso hace costoso refactorizar. No convertirlo en dogma; existen sistemas que necesitan separación temprana.

\subsection{Netflix}

Narrativa: microservicios compran independencia real de escala y despliegue en una organización preparada. El costo es plataforma y observabilidad.

\subsection{AWS Saga/Outbox}

Narrativa: los patrones organizan consistency gaps; no los hacen desaparecer. Leer también limitaciones oficiales.

\section{Demostraciones sugeridas}

\begin{uademo}
\textbf{Disponibilidad compuesta.} Calcular $0.999^4$ y luego agregar retries/latencia. Mostrar que cadena sincrónica cambia el sistema.
\end{uademo}

\begin{uademo}
\textbf{Dual write.} Dos estudiantes sostienen DB y broker; mover CRASH entre acciones.
\end{uademo}

\begin{uademo}
\textbf{Saga.} Tarjetas de pasos; el grupo decide retry, compensación o intervención manual para cada falla.
\end{uademo}

\section{Preguntas socráticas}

\begin{itemize}
\item ¿Qué autonomía compra esta frontera?
\item ¿Qué invariante deja de ser local?
\item ¿Quién es dueño del dato?
\item ¿La interacción necesita respuesta ahora?
\item ¿Qué pasa si el evento no se publica?
\item ¿La compensación puede deshacer realmente el efecto?
\item ¿Qué equipo recibe el pager a las 3 a.m.?
\item ¿Qué métrica demostraría que la extracción funcionó?
\end{itemize}

\section{Errores previsibles y corrección}

\begin{tabular}{p{0.38\textwidth}p{0.52\textwidth}}
\toprule
Error & Intervención \\
\midrule
“microservicios escalan” & Preguntar qué componente y qué carga. \\
“monolito = malo” & Mostrar monolito modular y transacción local. \\
“database per service = servidor por servicio” & Volver a ownership lógico. \\
“eventos desacoplan todo” & Preguntar por schema, orden y efecto. \\
“outbox da exactly-once” & Introducir relay duplicado. \\
“compensación = rollback” & Usar email o envío físico. \\
“orquestador = single point automáticamente” & Discutir persistencia/replicación y centralización lógica. \\
\bottomrule
\end{tabular}

\section{Criterios de evaluación formativa}

Una respuesta lograda incluye:
\begin{enumerate}
\item presión concreta para separar;
\item ownership e invariante;
\item comunicación y semántica;
\item nueva falla y resiliencia;
\item capacidad operativa;
\item métrica/condición de revisión.
\end{enumerate}

\section{Cómo mejorar la experiencia del estudiante}

\begin{itemize}
\item Use un caso cercano y manténgalo toda la clase.
\item Presente alternativas plausibles, no una respuesta “correcta” única.
\item Dé lenguaje para debatir sin slogans.
\item Incluya mini decisiones cada 15 minutos.
\item Reconozca explícitamente cuando un monolito es mejor.
\item Termine con una plantilla que puedan usar en el TP.
\end{itemize}

Esto aumenta percepción de respeto intelectual, aplicabilidad y claridad, factores centrales en una buena encuesta docente.

\section{Conexión con el Trabajo Final Integrador}

Cada grupo debe presentar una decisión de frontera con: capacidad, owner, dato, interacción, falla, operación y métrica. Retroalimente si la separación compra autonomía o solo cambia llamadas locales por red.

\section{Cierre y puente}

Cierre: “Al distribuir la arquitectura, también distribuimos la evidencia. La próxima clase estudia cómo reconstruir una request, una degradación y una causa a través de muchos servicios”.

\end{document}
